\section{引言：从他人经验中学习}

本节课的主题是``Everything you didn't want to know about LM architecture and training''——深入探讨语言模型架构和训练中那些看似琐碎但实际至关重要的细节。Tatsu Hashimoto 通过系统地调研过去一年发布的 19 个密集模型（dense model），以``趋同进化''的视角揭示了现代 Transformer 架构的设计共识与分歧。

\begin{importantbox}{本课核心主题}
最好的学习方式是亲手实践；次好的方式是从他人的经验中学习。我们无法训练所有这些 Transformer，但可以通过分析大量已发布模型的架构选择，提炼出可靠的设计原则。
\end{importantbox}

\begin{figure}[H]
\centering
\includegraphics[width=0.95\textwidth]{slides-images/slide-002.jpg}
\caption{课程大纲与目标：快速回顾标准 Transformer、分析大型语言模型的共同选择、探讨常见的架构变体\protect\footnotemark}
\end{figure}
\footnotetext{Slides 第3页。}

课程分为四个主要部分：

\begin{enumerate}
    \item \textbf{架构变体}（Architecture Variations）：归一化方式、激活函数、位置编码等
    \item \textbf{超参数选择}（Hyperparameters）：FFN 维度比、注意力头数、宽深比等
    \item \textbf{训练稳定性技巧}（Stability Tricks）：Z-loss、QK-norm 等
    \item \textbf{注意力变体}（Attention Variants）：GQA、MQA、滑动窗口注意力等
\end{enumerate}

\subsection{从原始 Transformer 到现代变体}

\begin{figure}[H]
\centering
\includegraphics[width=0.95\textwidth]{slides-images/slide-003.jpg}
\caption{原始 Transformer（Vaswani et al., 2017）的设计选择：正弦位置编码、ReLU FFN、Post-norm LayerNorm\protect\footnotemark}
\end{figure}
\footnotetext{Slides 第4页。}

原始 Transformer（2017）的关键设计选择：
\begin{itemize}
    \item \textbf{位置编码}：正弦/余弦函数
    \item \textbf{前馈层}：ReLU 激活，$\text{FFN}(x) = \max(0, xW_1 + b_1)W_2 + b_2$
    \item \textbf{归一化}：Post-norm（LayerNorm 在残差连接\textbf{之后}）
    \item \textbf{偏置项}：所有线性层和归一化层都有偏置
\end{itemize}

\begin{figure}[H]
\centering
\includegraphics[width=0.95\textwidth]{slides-images/slide-004.jpg}
\caption{课程作业实现的现代 Transformer 变体：Pre-norm、RoPE、SwiGLU、无偏置项\protect\footnotemark}
\end{figure}
\footnotetext{Slides 第5页。}

而现代 Transformer（如 LLaMA 系列）的关键变化：
\begin{itemize}
    \item \textbf{位置编码}：旋转位置编码（RoPE）
    \item \textbf{前馈层}：SwiGLU 激活
    \item \textbf{归一化}：Pre-norm（LayerNorm 在残差分支\textbf{之前}）
    \item \textbf{偏置项}：全部移除
\end{itemize}

\begin{knowledgebox}{19 个模型的``趋同进化''}
Tatsu 汇编了从 2017 年原始 Transformer 到 2025 年最新模型（Command A、Gemma 3、Qwen 2.5 等）的架构对比表。通过观察这些模型的选择趋势，我们可以清晰地看到架构设计的``趋同进化''——在某些维度上，所有人都收敛到了相同的选择。
\end{knowledgebox}

\begin{figure}[H]
\centering
\includegraphics[width=0.95\textwidth]{slides-images/slide-005.jpg}
\caption{19 个密集模型的架构选择汇总表：颜色编码展示了各维度上的趋同与分歧\protect\footnotemark}
\end{figure}
\footnotetext{Slides 第6页。}

\subsection{本章小结}

本课的核心方法论是``数据驱动的架构理解''：通过系统调研大量已发布模型的设计选择，识别出哪些是共识（所有人都做的），哪些是有益但非必要的（大多数人做的），哪些仍存在分歧（不同团队做不同选择）。

%% ========== Part 2: 架构变体 ==========

